
%%%%%%%%%%  scan idx 144  |  folio 133  %%%%%%%%%%


S G A 7


\underline{E X P O S I C I Ó N \ VII}


\underline{\underline{BIEXTENSIONES DE HACES DE GRUPOS}}


\underline{por A. Grothendieck}


0. \underline{Introducción}

La noción de \underline{biextensión} fue introducida por D. MUMFORD
[6], en el contexto de los grupos formales, para expresar cómodamente
las relaciones entre los grupos formales asociados a un esquema
abeliano y el esquema abeliano dual. En la presente exposición hacemos
un estudio sistemático de esta noción en el contexto general de los
haces de grupos sobre un sitio dado, y mostramos en particular que esta
noción se integra de manera notablemente sencilla en el formalismo
cohomológico clásico de $\overset{L}{\otimes}$ y
$\mathrm{R}\underline{\mathrm{Hom}}$ (3.6.4 y 3.6.5). En la exposición siguiente
explicitaremos los casos particulares más interesantes, y en especial
el de las biextensiones de esquemas en grupos por el grupo multiplicativo
$\underline{G}_{m}$ (trabajando con sitios tales como el sitio de
Zariski, el sitio fpqc o algún sitio intermedio); veremos en particular
que, si $A$ y $B$ son dos esquemas abelianos sobre un esquema $S$, las
biextensiones de $(A,B)$ por $\underline{G}_{m}$ poseen una
interpretación notable como, esencialmente, las correspondencias
divisoriales clásicas sobre $A,B$. Queda así explicitado de manera muy
cómoda el significado cohomológico de la noción de correspondencia
divisorial (ya implícito en el teorema de dualidad de Tate para
variedades abelianas sobre cuerpos $p$-ádicos). Esta interpretación de
las correspondencias divisoriales también resulta cómoda para «seguir»
el destino de una correspondencia de esta índole cuando se hacen
«degenerar»
%%%%%%%%%%  scan idx 145  |  folio 134  %%%%%%%%%%
$A$ y $B$ (supuestos definidos sobre el cuerpo de fracciones de un
anillo de valoración discreta) mediante la teoría de Néron: los
resultados de VIII nos mostrarán que se obtiene, por ejemplo, una
biextensión del par de componentes neutras de los modelos de Néron, y
por tanto una biextensión de las fibras especiales conexas de los
modelos de Néron, que desempeña el papel de una especialización de la
correspondencia divisorial inicial. Estos resultados se utilizarán en
IX, junto con el teorema de monodromía (I 3), para demostrar ciertos
resultados sobre los modelos de Néron de las variedades abelianas (el
más importante de ellos es el «teorema de reducción semiestable» para
el modelo de Néron).

En todo lo que sigue trabajaremos (salvo indicación expresa en
contrario) en un $\underline{U}$-topos fijo, denotado por
$\underline{T}$; es decir (SGA 4 IV), $\underline{T}$ es una categoría
equivalente a la categoría de los haces sobre un $\underline{U}$-sitio
adecuado (que puede tomarse, si se desea, igual a $\underline{T}$,
dotado de su topología canónica). A menudo identificaremos, en la
terminología, los objetos de $\underline{T}$ con los haces (de
conjuntos) sobre $\underline{T}$. Los grupos, extensiones de grupos,
etc., que consideraremos serán todos relativos a este topos
$\underline{T}$.


1. \underline{Complementos sobre las extensiones de grupos}

1.1.   Sean $P$ y $G$ dos grupos de $\underline{T}$, y

\begin{equation}\tag{1.1.1}
\begin{tikzcd}[column sep=large]
1 \arrow[r] & G \arrow[r, "i"] & E \arrow[r, "j"] & P \arrow[r] & 1
\end{tikzcd}
\end{equation}
% ESTRUCTURA: Una sola fila horizontal, 5 nodos, 4 flechas, todas flechas simples continuas hacia la
%   derecha (sin ganchos, sin cabezas dobles, sin trazos discontinuos, sin etiquetas de isomorfismo,
%   sin flechas verticales ni diagonales). NODOS (de izquierda a derecha): (1) «1»; (2) «G»; (3) «E»;
%   (4) «P»; (5) «1». FLECHAS: A1: nodo1 «1» --> nodo2 «G», hacia la derecha, flecha continua simple
%   de una cabeza, SIN etiqueta. A2: nodo2 «G» --> nodo3 «E», hacia la derecha, flecha continua simple
%   de una cabeza, etiqueta «i» ENCIMA del cuerpo de la flecha. A3: nodo3 «E» --> nodo4 «P», hacia la
%   derecha, flecha continua simple de una cabeza, etiqueta «j» ENCIMA del cuerpo de la flecha. A4:
%   nodo4 «P» --> nodo5 «1», hacia la derecha, flecha continua simple de una cabeza, SIN etiqueta. La
%   etiqueta «(1.1.1)» aparece en el margen IZQUIERDO de la línea, bien separada del diagrama, situado
%   hacia el centro del bloque de texto.
una extensión de $P$ por $G$; por tanto, $E$ es un grupo,
$j : E \longrightarrow P$ es un epimorfismo de grupos, e $i$ induce
un isomorfismo de $G$ con Ker $j$. Vamos a reinterpretar el dato de una
extensión de esta índole en términos diferentes, que nos resultarán
cómodos para definir y estudiar las biextensiones.


%%%%%%%%%%  scan idx 146  |  folio 135  %%%%%%%%%%

% Canon: SGA-SRC-001252.
1.1.2.  En primer lugar, el homomorfismo $i$ permite hacer actuar $G$
de dos maneras sobre $E$, a izquierda y a derecha, poniendo

\begin{equation*}
gxg' = i(g)xi(g')
\end{equation*}
para $g,g' \in G(S)$, $x \in E(S)$ ($S$ un objeto cualquiera de
$\underline{T}$). Si, gracias a $j$, se considera $E$ como un objeto
«sobre $P$», es decir, como un objeto del topos inducido
$\sisi{E_{/P}}{\underline T_{/P}}$, esto equivale a decir que el grupo
inducido $\sisi{G_P = G\times P}{G_P=G\times P}$ por $G$ en
$\sisi{\underline{T}_{/P}}{\underline T_{/P}}$ actúa sobre $E$ a
izquierda y a derecha. Además,
es bien sabido que una y otra acción hacen de $E$ un \underline{torsor}
(a izquierda y a derecha, respectivamente) sobre $P$, de grupo
$\sisi{G_P}{G_P}$,
y que se obtiene incluso, más precisamente, un \underline{bitorsor}
bajo $(\sisi{G_P}{G_P},\sisi{G_P}{G_P})$. Recordemos [2] que esto
significa que $E$ es un torsor a derecha $E_d$ para la acción derecha de
$\sisi{G_P}{G_P}$, y que la acción izquierda induce un isomorfismo de
$\sisi{G_P}{G_P}$ con el grupo de automorfismos
de $E_d$ (= haz de grupos de los automorfismos de $E$ que conmutan con
la acción derecha de $\sisi{G_P}{G_P}$ sobre $E$). Esta condición
equivale a la condición simétrica, a saber: la acción izquierda de
$\sisi{G_P}{G_P}$ sobre $E$ hace de $E$ un torsor a izquierda $E_s$, y
la acción derecha induce un isomorfismo del grupo opuesto
$\sisi{G_P^{\;o}}{G_P^{\;o}}$ de $\sisi{G_P}{G_P}$ con el grupo de
automorfismos de $E_s$.

1.1.2.1.  Cuando $\underline{T}$ es el «topos puntual», es decir, la
categoría de los conjuntos (equivalente a la categoría de los haces
sobre el espacio topológico puntual), de modo que $P$ y $G$ son grupos
ordinarios, el dato de un bitorsor $E$ bajo
$(\sisi{G_P}{G_P},\sisi{G_P}{G_P})$ equivale
manifiestamente al dato, para todo $p \in P$, de un bitorsor $E_p$ bajo
$(G,G)$ (a saber, la fibra de $E$ en $p$). Puede darse una
interpretación análoga en el caso general, observando que el dato de un
bitorsor $E$ bajo $(\sisi{G_P}{G_P},\sisi{G_P}{G_P})$ equivale al dato,
para todo «punto»
$p \in P(S)$ ($S$ un objeto cualquiera de $\underline{T}$), de un
bitorsor $E_p$ bajo
%%%%%%%%%%  scan idx 147  |  folio 136  %%%%%%%%%%
el par de grupos $(G_S,G_S)$ sobre $S$, de modo que la formación de
$E_p$ sea «compatible con los cambios de base para todo morfismo
$S' \longrightarrow S$»; es decir, si $p' : S' \longrightarrow P$ es
el morfismo compuesto, se supone dado además un isomorfismo de
% Canon: SGA-SRC-001253.
bitorsors de $\sisi{P_p}{E_{p'}}$ con la imagen inversa de
$\sisi{P_p}{E_p}$ por $S' \longrightarrow S$, estando estos
isomorfismos sometidos a la condición usual de transitividad para el
caso de un compuesto de dos morfismos
$S'' \longrightarrow S' \longrightarrow S$. Al explicitar este
enunciado en lenguaje canónico (SGA 1 VI), se reduce a la tautología
siguiente: el dato de un bitorsor bajo el par de grupos
$(\sisi{G_P}{G_P},\sisi{G_P}{G_P})$
sobre $P$ equivale a una sección cartesiana, sobre la categoría
$\underline{T}_{/P}$, de la categoría fibrada de los bitorsors de grupo
$(G_S,G_S)$ sobre una base variable
$S \in \mathrm{ob}\ \underline{T}_{/P}$; más precisamente aún, el
funtor «valor en $P$» induce una equivalencia de la categoría de las
secciones cartesianas en cuestión con la categoría de los bitorsors bajo
$(\sisi{G_P}{G_P},\sisi{G_P}{G_P})$. Este enunciado resulta, por
supuesto, trivialmente del hecho
de que $P$ es objeto final de la categoría base $\underline{T}_{/P}$ de
nuestra categoría fibrada; por tanto, para \underline{toda} categoría
fibrada sobre $\underline{T}_{/P}$, el funtor «valor en $S$» es una
equivalencia entre la categoría de las secciones cartesianas de dicha
categoría fibrada y la categoría fibra en $P$. No obstante, pese a su
naturaleza tautológica, esta interpretación nos será útil para expresar
ciertas compatibilidades en un lenguaje próximo al caso «conjuntista»,
utilizando puntos con valores en un objeto $S$ en lugar de puntos
ordinarios.
